\cotefile{157-3}
\scanpage{114}
% line head | text | 63
\mspage{63}
% line 1 | text | donne ceci
\unsure{donne} ceci :

% line 2–10 | text | Proposition 2
\uline{Proposition 2} : Considérons une situation de recollement
de deux topos $\mathscr U$, $\mathscr T$ par $\varphi:\operatorname{Fais}\mathscr U\unsure{-|\to}\operatorname{Fais}\mathscr T$,
foncteur exact à gauche. Supposons que $\mathscr U=\operatorname{Top}(U)$,
$\mathscr T=\operatorname{Top}(T)$, avec $U$ et $T$ dans $\mathrm{Cat}$. Pour que
le topos recollé $\unsure{\mathscr X|\mathscr Z}$ soit équivalent à un topos
$\operatorname{Top}(X)$, avec $X$ dans $\mathrm{Cat}$, il f. et s. que
le foncteur de recollement $\varphi$ commute aux
$\lim_{\leftarrow}$ \unsure{(quelc.),|(quelc.,)} ou ce qui revient au $=$,
qu'il ait un adjoint à gauche :
% line 11 | display | ψ, Fais T, Fais U, the two presheaf categories
\[
\psi:\underset{\substack{\shortparallel\\T^\wedge}}{\underline{\operatorname{Fais}\mathscr T}}
\longrightarrow\underset{\substack{\shortparallel\\U^\wedge}}{\underline{\operatorname{Fais}\mathscr U}} .
\]

% line 12–16 | text | On s'attend dès lors
On s'attend dès lors que \deleted{\unsure{cette}} la donnée de
recollement $\varphi$, ou ce qui revient au $=$ le
foncteur $\psi$ \unsure{a}\deleted{\unsure{yant}} un adjoint à droite, i.e.
commutant aux $\lim_{\rightarrow}$ quelc., revienne à
celle d'une donnée de recollement $\unsure{H|M}$
% line 17 | display | H, U with a small circle, T, Ens
\[
H:U^{\mathrm o}\times T\longrightarrow\mathrm{Ens}
\]
% line 18–24 | text | pour U et T
pour $U$ et $T$. Mais on sait bien que
les foncteurs $\psi:T^\wedge\longrightarrow\mathscr C$ dans une
catégorie \inserted{$\mathscr C$} avec $\lim_{\rightarrow}$ quelc., qui commutent
aux dites $\lim_{\rightarrow}$, forment une catégorie
équivalente (\unsure{par|pour} le foncteur $\psi\longmapsto\psi\mid T\,(\subset T^\wedge)$)
à celle des \deleted{\unsure{préf.}} foncteurs
% line 25 | display | ψ₀, T, C, U with wedge, Hom
\[
\psi_0:T\longrightarrow\mathscr C
\quad \bigl(=\underset{\substack{\shortparallel\\\underline{\operatorname{Hom}}(\deleted{\unsure{U}}U^{\mathrm o},\mathrm{Ens})}}{U^\wedge}
\text{ en l'occurrence}\bigr).
\]
% archive sheet number: 114
